\section{Literature Review}
\label{sec:literature}

This paper connects three strands of the literature: labor market scarring, informality in developing countries, and the differential impacts of COVID-19 by occupation.

\subsection{Labor Market Scarring}

The scarring hypothesis holds that adverse labor market shocks can produce persistent effects on workers' employment quality, wages, and formality status long after the initial shock dissipates. \citet{arulampalam2001} document that unemployment spells in the UK carry a wage penalty of 6\% that persists for over a decade. In developing countries, \citet{cruces2012} show that exposure to economic crises in Argentina durably reduced formal sector attachment. The mechanism is that workers displaced into informality face barriers to re-entry---skill depreciation, loss of employer networks, and stigma---creating a ``trap'' that outlasts the shock.

\subsection{Informality in Peru}

Peru has one of the highest informality rates in Latin America, with over 70\% of workers lacking formal employment relationships. \citet{bosch2025} study the effect of enforcement letters on firm formalization, finding modest but significant effects on social security enrollment among large firms. However, the worker-side experience of informality transitions remains underexplored. The COVID-19 pandemic provides a quasi-natural experiment: the lockdowns forced millions of formal workers into informal arrangements, creating a large cohort of ``involuntary informals'' whose subsequent trajectory can be tracked.

\subsection{COVID-19 and Occupation-Based Exposure}

\citet{dingel2020} introduced the teleworkability index, classifying occupations by the feasibility of remote work based on O*NET task descriptions. This index has been widely used as a source of identifying variation for COVID-19's differential labor market impacts. \citet{saltiel2020} adapted the index for developing countries, finding that only 13\% of workers in low- and middle-income countries can work from home, compared to 37\% in the U.S. For Peru specifically, our data shows that 74.6\% of employed workers are in contact-intensive occupations (teleworkability score $< 0.20$), making the country an important setting for studying the long-run consequences of occupation-based COVID exposure.
